% !TEX root = ../thesis.tex

\chapter{Methodology}
\label{ch:methodology}

This chapter consolidates the experimental methodology of the thesis. It
first describes the benchmark dataset and the preprocessing pipeline that
produces a consistent input table for all generators
(Section~\ref{sec:method:data}), then specifies the four synthetic-data
generation models compared in this work
(Section~\ref{sec:method:generators}), and finally defines the two-pillar
evaluation framework used to score every generator
(Section~\ref{sec:method:evaluation}).

\section{Dataset and Preprocessing}
\label{sec:method:data}

\subsection{The Diabetes 130-US Hospitals Dataset}
\label{sec:data:dataset}

The benchmark used throughout this thesis is the \emph{UCI Diabetes 130-US
Hospitals} dataset~\cite{diabetes_db,strack2014impact}. It contains 101{,}766
inpatient encounters from diabetic patients treated across 130 U.S.\ hospitals
between 1999 and 2008. Each record describes a hospital admission of 1--14
days and includes 47 mixed-type features covering demographics (age, gender,
race), admission characteristics, laboratory test results, diagnoses,
medication use, and prior inpatient, outpatient, and emergency visits. The
data consists primarily of categorical and integer variables and represents
detailed clinical, administrative, and treatment-related information collected
during the patient's hospital stay. Missing values are present, and no
predefined data splits are provided, enabling flexible
train--validation--test partitioning. The binary target used throughout this
work is \texttt{readmitted}: hospital readmissions of any kind are encoded as
the positive class and ``no readmission'' as the negative class.

\paragraph{Scope.} The proposal contemplated MIMIC-III as a secondary
benchmark; the final thesis is restricted to the Diabetes 130-US Hospitals
dataset to keep training budgets and evaluation runs manageable across four
generators and seventeen metrics. Cross-dataset generalisation is discussed as
future work in Section~\ref{sec:conclusion:future}.

\subsection{Exploratory Data Analysis}
\label{sec:data:eda}

The raw CSV contains 50 columns: 13 integer columns and 37 object columns
when loaded directly with \texttt{pandas}. The integer fields include
encounter and patient identifiers, utilisation counts, procedure counts,
medication counts, diagnosis counts, and length of stay. Most remaining
fields are categorical, including demographic variables, admission and
discharge codes, ICD-9 diagnosis codes, laboratory-result indicators,
medication-change indicators, and the readmission target. The preprocessing
loader treats \texttt{?} as a missing value and then casts known count and
identifier fields to numeric type while leaving the remaining columns as
categorical variables. Missingness is concentrated in a small number of
clinical and administrative fields: weight (96.86\% missing),
max-glucose-serum (94.75\%), \texttt{A1Cresult} (83.28\%), medical specialty
(49.08\%), and payer code (39.56\%), with lower missingness in \texttt{race}
(2.23\%), \texttt{diag\_3} (1.40\%), \texttt{diag\_2} (0.35\%), and
\texttt{diag\_1} (0.02\%). The pipeline removes \texttt{weight},
\texttt{payer\_code}, and \texttt{medical\_specialty} because their
missingness is high and their values are not needed for the downstream
synthetic-data comparison; \texttt{race} is also removed after resolving
multiple encounters, avoiding use of a sensitive demographic field, while
\texttt{A1Cresult} and the diagnosis fields are retained because their
missingness can be represented or imputed after clinically motivated
encoding.

Several medication variables are highly imbalanced in the raw dataset
because almost all encounters record \texttt{No}. Fifteen medication fields
have fewer than 1\% non-mode values, including rare drugs such as
\texttt{examide}, \texttt{citoglipton}, \texttt{troglitazone}, and several
combination medications. The pipeline keeps these medication fields, but
converts each from the original \texttt{No}, \texttt{Steady}, \texttt{Up},
and \texttt{Down} representation into a binary exposure indicator. This
preserves whether a drug was used while reducing sparse categorical levels
that would be difficult for tabular generators to model reliably.

The raw target has three categories: \texttt{NO}, \texttt{>30}, and
\texttt{<30}. In the full raw dataset, \texttt{NO} accounts for 54{,}864
encounters (53.91\%), \texttt{>30} for 35{,}545 encounters (34.93\%), and
\texttt{<30} for 11{,}357 encounters (11.16\%). For this thesis, the two
readmission categories are collapsed into a single positive class because
the main modelling task is whether any readmission occurred. After
multiple-encounter resolution and binary target encoding, the processed
dataset contains 47{,}188 non-readmitted records (65.98\%) and 24{,}330
readmitted records (34.02\%). The stratified training split preserves this
imbalance with 37{,}750 negatives and 19{,}464 positives. This class
distribution is used as a baseline property that synthetic generators
should reproduce.

\subsection{Feature Engineering Pipeline}
\label{sec:data:fe}

The preprocessing and feature-engineering code is implemented in the
feature-engineering module and executed by the data-preprocessing script.
The script loads the raw data, applies the engineering pipeline, writes the
full preprocessed table, and then creates stratified train and test CSV
files. The transformation order is fixed so that row-level deduplication and
clinically meaningful grouping occur before imputation and label encoding.
The raw dataset contains repeated encounters for some patients: although
there are 101{,}766 encounter rows, there are 71{,}518 unique values of
\texttt{patient\_nbr}; 16{,}773 patients appear more than once, and the
maximum number of encounters for a single patient is 40. To avoid
patient-level leakage between train and test splits,
\texttt{process\_multiple\_encounters} keeps one row per patient by sorting
on \texttt{patient\_nbr} and \texttt{time\_in\_hospital}, retaining the
encounter with the longest hospital stay, and dropping duplicate patient
rows. \texttt{patient\_nbr} is then removed because it is a direct patient
identifier, and \texttt{race} is also removed to reduce the use of
sensitive demographic information.

\label{sec:data:fe:age}%
\texttt{consolidate\_age} converts the original decade-based age buckets
into ordered numeric values: the final experiments use a 1--10 ordinal
encoding where \texttt{[0-10)} maps to 1, \texttt{[10-20)} to 2, and so on
through \texttt{[90-100)}, which maps to 10. This preserves the natural
ordering of age while avoiding a high-cardinality string representation. An
earlier preprocessing variant collapsed the ten decades into a three-level
grouping (young / middle-aged / elderly) to reduce categorical sparsity,
but a direct ablation against the ten-bin encoding showed that the coarser
grouping substantially degraded downstream synthetic-data utility for the
LLM-based generator while leaving the GAN baselines largely unaffected. The
native ten-bin encoding is therefore retained for all experiments in this
thesis. Quantitative results of the ablation are reported in
Section~\ref{sec:results:rq2:utility:age-binning}. Three sparse fields ---
weight, payer code, and medical specialty --- are also dropped to reduce
dimensionality and prevent the synthetic-data generators from spending
capacity on poorly observed variables.

\texttt{change\_diag\_columns} maps \texttt{diag\_1}, \texttt{diag\_2}, and
\texttt{diag\_3} from raw ICD-9 strings into broader diagnosis groups using
\texttt{\_icd9\_to\_group}: codes beginning with 250 are mapped to
\texttt{Diabetes}, and numeric ICD-9 ranges are grouped as circulatory,
respiratory, digestive, genitourinary, neoplasm, musculoskeletal, and injury
categories. External cause and supplementary codes beginning with \texttt{E}
or \texttt{V}, missing values, and unrecognised codes are mapped to
\texttt{Other}. This step reduces thousands of possible diagnosis-code
values into nine clinically interpretable categories, which are converted to
integer labels later by the label-encoding step.
\texttt{change\_medication\_columns} transforms 23 medication columns into
binary indicators: for each medication, \texttt{No} is encoded as 0 and
every other value, including \texttt{Steady}, \texttt{Up}, and \texttt{Down},
is encoded as 1. The encoded set includes individual diabetes medications
and combination therapies such as metformin, insulin, sulfonylureas,
thiazolidinediones, and metformin-based combinations. The binary
representation emphasises whether a medication was prescribed or changed
during the encounter, rather than attempting to model sparse dosage
direction categories separately.

\texttt{encode\_a1cresult} converts \texttt{A1Cresult} into an ordinal
laboratory indicator: missing A1C values are encoded as 0, \texttt{Norm} as
1, \texttt{>7} as 2, and \texttt{>8} as 3, keeping the absence of an A1C
test as an explicit value while preserving the severity ordering of
observed test results. \texttt{impute\_data} is then applied: categorical
columns are imputed with the most frequent observed value, while numeric
columns are imputed with the median, both via \texttt{SimpleImputer}. After
imputation, the preprocessed dataset contains no missing values.
\texttt{label\_encode} converts all remaining object-type feature columns
into integer labels using \texttt{LabelEncoder}, excluding the target.
Finally, \texttt{target\_encode} maps the original three-class readmission
variable into a binary outcome: \texttt{NO} is encoded as 0, while both
\texttt{>30} and \texttt{<30} are encoded as 1. The binary formulation
treats any hospital readmission as the actionable positive outcome and is
used consistently for model training, synthetic-data evaluation, and
downstream metric computation.

\subsection{Train/Test Split}
\label{sec:data:split}

The preprocessing script writes three CSV files under \texttt{thesis/data}:
the full preprocessed dataset, the training split, and the test split. The
full processed dataset contains 71{,}518 rows and 45 columns. The split is
created with \texttt{train\_test\_split}, \texttt{test\_size = 0.2}, and
\texttt{random\_state = 42}. The target column is passed as the
stratification variable, preserving the binary readmission ratio across
splits.

The final training set contains 57{,}214 rows and 45 columns, with 37{,}750
negative records and 19{,}464 positive records. The held-out test set
contains 14{,}304 rows and 45 columns, with 9{,}438 negative records and
4{,}866 positive records. Both splits contain no missing values and no
object-type columns after preprocessing.

\section{Synthetic Data Generation Methods}
\label{sec:method:generators}

% Expansion of proposal §5.2 (Proposed Models). One subsection per generator
% with architecture, hyperparameters, training protocol, and (where
% applicable) the differential-privacy mechanism.

\subsection{Selection Rationale}
\label{sec:methods:rationale}

% Migrated from proposal §5.2 (selection rationale).
The four generators evaluated in this thesis are chosen to provide a
$2 \times 2$ coverage of the design space: $\{\text{GAN},\text{LLM}\}$ on the
modelling axis, and $\{\text{non-DP},\text{DP}\}$ on the privacy axis. Within
each axis the choice is grounded in prior work:

\begin{itemize}\itemsep0.25em
  \item \textbf{HealthGAN (non-DP GAN baseline).} HealthGAN was designed
        specifically for mixed-type EHR tables, comes with a complete
        secure training-to-export workflow, and introduces nearest-neighbour
        Adversarial Accuracy (AA) to co-measure resemblance and privacy
        (TrainAA/TestAA and privacy loss), yielding competitive TSTR utility
        while keeping a small footprint~\cite{healthgan,yale2020synthesizing,dash2020medical}.
  \item \textbf{PATE-GAN (DP GAN).} PATE-GAN enforces formal
        $(\varepsilon,\delta)$-DP by decoupling privacy from adversarial
        training: multiple teacher discriminators are trained on disjoint
        partitions, their votes are noisily aggregated, and a student
        discriminator learns from these noisy labels so the generator inherits
        DP by post-processing~\cite{pategan,Papernot2017PATE}.
  \item \textbf{Pythia LLM (non-DP LLM baseline).} The benchmark study by
        \textcite{Miletic2024} evaluates the Pythia family of decoder-only
        LLMs for tabular health data and shows they can match or exceed
        established tabular GANs. The authors introduce \emph{Tabula}, a
        framework that treats each tabular row as a text sequence. Their
        evaluation focuses on utility rather than privacy.
  \item \textbf{Pythia-DP (DP LLM, scope extension).} Beyond the originally
        proposed three models, this thesis adds a differentially private
        variant of Pythia trained with DP-SGD via
        Opacus~\cite{Yousefpour2021Opacus,Abadi2016DPSGD}, enabling a
        within-family DP-cost comparison parallel to the
        HealthGAN $\to$ PATE-GAN comparison.
\end{itemize}

\subsection{HealthGAN: WGAN-GP Baseline}
\label{sec:methods:healthgan}

% Code references:
%   - SDV encoding/decoding: ../../healthgan/generators/sdv_converter.py
%   - WGAN-GP training:      ../../healthgan/generators/wgan_for_mac.py

\subsubsection{SDV Encoding and Decoding}
\label{sec:methods:healthgan:sdv}

% TODO: Describe \texttt{encode}, \texttt{decode}, and the snap-to-nearest
% post-processing in sdv_converter.py.
%  - Categorical columns mapped via truncated normals (categorical())
%  - Ordinal columns mapped via additive smoothing + truncated normals
%  - Binary columns mapped via truncated beta (when --beta flag set)
%  - Numeric columns min-max normalised to [0, 1]
%  - Decoding inverts each, with a snap-to-nearest pass for low-cardinality
%    numeric columns (e.g. age) that should not be reconstructed continuously.
% Include the FORCE_ORDINAL set ({"age"}) and the rationale.

\subsubsection{WGAN-GP Architecture}
\label{sec:methods:healthgan:arch}

% Architecture (from wgan_for_mac.py):
%  Generator: input ~ N(0, I)^{100} → Dense(2F, ReLU) → Dense(1.5F, ReLU) → Dense(F, sigmoid)
%  Critic:    input(F) → Dense(64, LeakyReLU) → Dense(128, LeakyReLU) → Dense(256, LeakyReLU) → Dense(1)
%  F = 45 features after preprocessing.

\subsubsection{Training Loss: Wasserstein + Gradient Penalty}
\label{sec:methods:healthgan:loss}

% TODO: WGAN loss + gradient penalty term \cite{Gulrajani2017WGANGP},
% lambda = 10. Brief equation block.

\subsubsection{Hyperparameters and Training Configuration}
\label{sec:methods:healthgan:hp}

% From wgan_for_mac.py params:
%  - num_epochs = 100,000
%  - critic_iters = 5
%  - base_nodes = 64
%  - samples_per_file = 57,214
%  - Adam(lr=1e-4, beta1=0.5, beta2=0.9) for both G and D

\subsubsection{Variant Selection: samples\_0 vs.\ samples\_3 vs.\ samples\_6}
\label{sec:methods:healthgan:variants}

% This subsubsection reports notebook cells 53–60: comparing the three decoded
% samples by KS + correlation difference, justifying which variant feeds the
% cross-method comparison in Chapter 6.
% TODO:
%  - Table: per-variant mean KS, mean |corr diff|, similarity rank.
%  - Figure: PCA overlap (variant deep-dive cell 58).
%  - State which variant was chosen and why.

\subsubsection{Rationale for Excluding a DP-Augmented HealthGAN Variant}
\label{sec:methods:healthgan:nodp}

% Migrated from thesis_conversation.md scope-decision note.
A natural question is whether differential privacy could be applied
directly to the HealthGAN training loop --- for instance by replacing
the Adam optimiser of the critic with a DP-SGD variant
\cite{Abadi2016DPSGD} via \texttt{tensorflow\_privacy}. This thesis does
not pursue that direction; the DP-GAN comparison is instead realised
through the architecturally distinct PATE-GAN model
(\S\ref{sec:methods:pategan}). Four technical considerations motivate
this choice.

First, the gradient-penalty term that distinguishes WGAN-GP from
weight-clipping WGAN \cite{Arjovsky2017WGAN, Gulrajani2017WGANGP} is
itself a function of real data, as it is computed over interpolated
samples drawn between real and generated points. Its gradient cannot be
cleanly clipped on a per-example basis as DP-SGD requires, so
accommodating DP-SGD would necessitate dropping the gradient penalty
and reverting to weight-clipping Wasserstein training, sacrificing the
architectural property that motivated the WGAN-GP variant in the first
place. Second, the training protocol applies five critic updates per
generator update; because the critic is the sole component that
directly processes real records, this $5{:}1$ ratio implies the privacy
budget would be consumed five times faster than in a configuration
performing one private gradient computation per logical batch, as is
the case in the Pythia-DP setting (\S\ref{sec:methods:pythia-dp}).
Third, the injection of Gaussian noise into critic gradients --- the
core operation of DP-SGD --- is well documented to exacerbate the mode
collapse and non-convergence that already characterise GAN training on
tabular health data, with empirical results in the DP-GAN literature
\cite{Xie2018DPGAN, pategan} demonstrating sharp utility cliffs at
budgets of $\varepsilon \leq 10$. Fourth, the HealthGAN implementation
relies on TensorFlow~1 compatibility mode with eager execution
disabled, and the per-example gradient computation required by
\texttt{tensorflow\_privacy} is substantially slower under graph
execution, rendering the $100{,}000$-epoch training schedule
computationally prohibitive within the resource envelope of this
project.

Beyond these technical obstacles, the experimental design already
provides a controlled DP-GAN baseline through PATE-GAN, which attains
formal $(\varepsilon,\delta)$-DP through a fundamentally different
mechanism --- noisy aggregation of teacher discriminators
\cite{pategan, Papernot2017PATE} --- rather than by injecting noise
into the critic's gradients. The HealthGAN $\to$ PATE-GAN pairing
therefore isolates the cost of formal DP on the GAN side without
confounding it with an unstable adaptation of an architecture that was
not designed for per-example gradient clipping. HealthGAN is retained
as the non-private GAN utility baseline.

\subsection{PATE-GAN: Differentially Private GAN}
\label{sec:methods:pategan}

% Code references:
%   - Pategan/pate_gan.py
%   - Pategan/generate_pategan_synthetic.py

\subsubsection{Teacher--Student Framework}
\label{sec:methods:pategan:teacher}

% TODO: Describe the k=10 teacher discriminators trained on disjoint shards;
% the student discriminator trained on noisily aggregated teacher labels;
% and the generator updated against the student.

\subsubsection{PATE Voting with the Laplace Mechanism}
\label{sec:methods:pategan:laplace}

% TODO: Equation for the noisy aggregation
%   y_hat = argmax_c (\sum_t 1[T_t(x) = c] + Lap(0, \lambda))
% with \lambda = 1.0 in the configuration used.

\subsubsection{Privacy Accounting via the Moments Accountant}
\label{sec:methods:pategan:accountant}

% TODO: Brief summary of the moments accountant
% \cite{Abadi2016DPSGD,Mironov2017RDP} as adapted by PATE-GAN \cite{pategan}.

\subsubsection{Architecture}
\label{sec:methods:pategan:arch}

% Generator: input(F) → Dense(4F, tanh) → Dense(4F, tanh) → Dense(F, sigmoid)
% Student:   input(F) → Dense(F, ReLU) → Dense(1)

\subsubsection{Privacy Budget and Hyperparameters}
\label{sec:methods:pategan:hp}

% k = 10 teachers, n_s = 1, batch_size = 64, lambda = 1.0
% Target privacy: epsilon = 5.0, delta = 1e-5.
% RMSprop(lr=1e-4) for both student and generator.
% Student weights clipped to [-0.01, 0.01].

\subsection{Pythia: Non-DP LLM-based Generation}
\label{sec:methods:pythia}

% Code references:
%   - pythia/pythia_tabular.py
%   - pythia/generate_pythia_synthetic.py

\subsubsection{Base Model}
\label{sec:methods:pythia:base}

% EleutherAI/pythia-70m \cite{Biderman2023Pythia}. Justify the small size
% choice (training cost / commodity GPU feasibility), referencing
% \cite{Miletic2024} who report Pythia-70m as a viable starting point.

\subsubsection{Tabular-to-Text Serialisation}
\label{sec:methods:pythia:serialisation}

% Format: "Class_{label} | col_name=value | col_name2=value2 | ..."
% Each row is converted to a single text sequence with the class token first
% so that conditional generation can prefix the class.

\subsubsection{Schema Derivation}
\label{sec:methods:pythia:schema}

% Per-column schema captures dtype, numeric bounds, discrete support, and
% category set (if applicable). Used downstream for parsing/coercion.

\subsubsection{LoRA Fine-tuning}
\label{sec:methods:pythia:lora}

% LoRA \cite{Hu2022LoRA} with r=8, alpha=16, dropout=0.05. Only LoRA
% parameters are trainable; the base Pythia model is frozen. Trainer:
% HuggingFace Trainer with DataCollatorForLanguageModeling (CLM objective).
%
% Hyperparameters (defaults):
%   epochs = 10
%   batch_size = 8
%   lr = 2e-5
%   max_length = 512

\subsubsection{Class-Conditioned Generation}
\label{sec:methods:pythia:generation}

% Sampling: temperature = 0.8, top_p = 0.95, max_new_tokens = 512.
% Prompt: "Class_{0|1} | col_name=" — let the model autoregressively complete
% the row.

\subsubsection{Schema-Aware Parsing, Coercion, and Retry}
\label{sec:methods:pythia:parsing}

% TODO:
%  - Regex extraction of key=value pairs.
%  - Coercion: range clip for numeric, discrete snapping, category fallback.
%  - Validation: ≥10 % of non-target columns must be present; otherwise retry.
%  - Fallback: resample from accepted rows; last resort schema-valid random rows.
%  - max_retries_per_row = 8.

\subsection{Pythia-DP: Differentially Private LLM}
\label{sec:methods:pythia-dp}

% Code references:
%   - pythia/pythia_tabular_dp.py
%   - pythia/generate_pythia_synthetic_dp.py

\subsubsection{DP-SGD via Opacus}
\label{sec:methods:pythia-dp:opacus}

% PrivacyEngine() from Opacus \cite{Yousefpour2021Opacus} wraps the model,
% optimizer, and dataloader. ModuleValidator.fix() is applied to make the
% model Opacus-compatible.

\subsubsection{Per-Example Gradient Clipping}
\label{sec:methods:pythia-dp:clipping}

% max_grad_norm = 1.0 — Opacus computes per-sample gradients (rather than
% per-batch) and clips each sample's gradient norm before averaging and
% adding noise.

\subsubsection{Effective Batch Size}
\label{sec:methods:pythia-dp:batch}

% per_device_batch_size = 32, gradient_accumulation_steps = 16 →
% effective_batch_size = 512. The privacy accountant tracks the effective
% batch size, not the per-device one.

\subsubsection{Privacy Budget and Achieved Epsilon}
\label{sec:methods:pythia-dp:epsilon}

% target_epsilon = 5.0, target_delta = 1e-5. Opacus computes the noise
% multiplier needed to reach target epsilon over the planned number of
% training steps. Achieved epsilon is reported in run_metadata_dp.json after
% training and should be cited verbatim in Chapter 6.

\subsubsection{Differences from the Non-DP Pipeline}
\label{sec:methods:pythia-dp:diff}

% TODO: Table or bulleted comparison highlighting:
%  - Trainer (HF) → manual training loop with PrivacyEngine
%  - Per-example gradients (memory cost)
%  - Effective batch size 512 vs. 8 in non-DP run
%  - Loss-curve dynamics typically shifted by noise injection

\subsection{Implementation Summary}
\label{sec:methods:summary}

% TODO: One synthesis table comparing all four generators along:
%  - Architecture
%  - Formal DP guarantee (yes/no, mechanism)
%  - Library stack (TensorFlow 1.x / PyTorch + Opacus / HF Transformers + LoRA)
%  - Training time on the available hardware
%  - Output file path under thesis/data/<generator>/

\section{Evaluation Framework}
\label{sec:method:evaluation}

% Expansion of proposal §5.3 (Evaluation Metrics). Methodology only — every
% definition, formula, and rationale lives here. Numerical results live in
% Chapter~\ref{ch:results}.

\subsection{Two Pillars: Utility and Privacy}
\label{sec:eval:pillars}

% TODO: One short paragraph framing the two-pillar structure adopted in this
% thesis, anchored in \cite{hernandez2022synthetic,Murtaza2023}.

This section defines a two-pillar evaluation framework: \emph{utility} and
\emph{privacy}. The utility pillar is treated as a single concept with two
complementary branches: \emph{statistical similarity} between the synthetic
and real datasets (Section~\ref{sec:eval:utility:similarity}) and
\emph{predictive performance} of downstream models trained on the synthetic
data (Section~\ref{sec:eval:utility:predictive}). The privacy pillar is
defined in Section~\ref{sec:eval:privacy} via instance-level distance and
inference-attack metrics. Section~\ref{sec:eval:aggregate} introduces the
aggregation views that combine both pillars into a single comparative
picture.

\subsection{Utility Metrics}
\label{sec:eval:utility}

% Maps to data_analysis.ipynb cells 29–41.

The utility pillar combines two complementary aspects. The first is whether
the synthetic dataset \emph{looks like} the real one in a distributional
sense; the second is whether the synthetic dataset \emph{acts like} the real
one when used as training data for downstream predictive and survival models.
The same set of synthetic records is evaluated under both views.

\subsubsection{Statistical Similarity to the Real Data}
\label{sec:eval:utility:similarity}

% Maps to data_analysis.ipynb cells 29–34. These metrics correspond to what
% \cite{Murtaza2023,hernandez2022synthetic} term resemblance / fidelity; in
% this thesis they are the first branch of the utility pillar.

\paragraph{Per-Column Distributional Match.}
\label{sec:eval:utility:similarity:dist}
% TODO: KDE for continuous, bar plots for categorical. Visual diagnostic only.

\paragraph{Summary Statistics (Mean, Standard Deviation).}
\label{sec:eval:utility:similarity:stats}
% TODO: Per-column mean / std table; relative differences from real.

\paragraph{Pairwise Correlation Difference (PCD).}
\label{sec:eval:utility:similarity:pcd}
% Following \cite{goncalves2020generation}, define
%   PCD = || corr(Real) - corr(Synth) ||_F
% the Frobenius norm of the Pearson correlation matrix difference.
% Smaller PCD indicates better preservation of joint linear structure.

\paragraph{Kolmogorov--Smirnov Two-Sample Test.}
\label{sec:eval:utility:similarity:ks}
% TODO:
%  - Definition of KS statistic.
%  - Per-column KS computed via scipy.stats.ks_2samp.
%  - Note: the proposal listed KL divergence; the implemented pipeline uses
%    KS as a non-parametric alternative requiring no PMF binning. Discuss
%    in §7.5.5 (Limitations) and Ch 2 (background) — both metrics share the
%    same purpose of distributional similarity.

\paragraph{PCA-Based Manifold Overlap.}
\label{sec:eval:utility:similarity:pca}
% Standardise (StandardScaler), fit PCA on combined Real + Synth, project to
% 2D, scatter plot per source. Visual assessment of manifold overlap.

\subsubsection{Predictive Performance of Downstream Models}
\label{sec:eval:utility:predictive}

% Maps to data_analysis.ipynb cells 35–41.

\paragraph{Train--Test Regimes: TRTR and TSTR.}
\label{sec:eval:utility:predictive:tstr}
% TRTR (train on real, test on real) baseline; TSTR (train on synthetic,
% test on real) per generator. The proposal additionally listed TRTS and
% TSTS; these are scope-narrowed in the implemented pipeline. Note this in
% §7.5.5 (Limitations).

\paragraph{Classifier Suite.}
\label{sec:eval:utility:predictive:classifiers}
% Three downstream classifiers:
%  - Logistic Regression (max_iter=1000)
%  - Random Forest (n_estimators=100)
%  - Gradient Boosting (n_estimators=100)

\paragraph{Performance Metrics.}
\label{sec:eval:utility:predictive:metrics}
% Per (Train Source × Classifier) cell, report:
%  - Accuracy
%  - F1 (zero_division=0)
%  - AUC-ROC
%  - Precision (zero_division=0)
%  - Recall (zero_division=0)
% AUC-ROC is the headline metric (class-imbalance robust); F1 secondary.

\paragraph{ROC Curve Analysis.}
\label{sec:eval:utility:predictive:roc}
% Gradient-Boosting ROC curves overlaid for Real (TRTR) vs. each synthetic
% source. Compare AUC and the false-positive trade-off.

\paragraph{Feature Importance Preservation.}
\label{sec:eval:utility:predictive:fi}
% Random Forest .feature_importances_ ranking per training source.
% Comparison metric: Spearman rank correlation of feature importances
% (Real vs. each Synth). Higher ρ = better preservation.

\paragraph{Cox Proportional Hazards Regression.}
\label{sec:eval:utility:predictive:cox}
% Cox model from lifelines.
%  - Duration: time_in_hospital
%  - Event:    readmitted
%  - Covariates: age, gender, num_medications, num_lab_procedures,
%    num_procedures, number_diagnoses, number_inpatient, number_emergency,
%    number_outpatient, insulin, diabetesMed, change, admission_type_id.
% Reports hazard ratios with 95 % confidence intervals.

\paragraph{Cox CI Overlap Analysis.}
\label{sec:eval:utility:predictive:cox-overlap}
% For each covariate, compute fraction of overlap between Real-fit CI and
% Synth-fit CI. Aggregate fraction (covariates with overlapping CI) is a
% scalar utility score: higher = better preservation of causal/risk
% relationships.

\subsubsection{Privacy Metrics}
\label{sec:eval:privacy}

% Maps to data_analysis.ipynb cells 42–48.

\paragraph{Nearest-Neighbour Distance.}
\label{sec:eval:privacy:nn}
% For each synthetic record, compute the L2 distance (in StandardScaler
% feature space) to its nearest real record. Report median NN distance per
% generator. Larger = better privacy (no near-duplicates of real patients).

\paragraph{Nearest-Neighbour Distance Ratio (NNDR).}
\label{sec:eval:privacy:nndr}
% Define
%   NNDR = median(d_{Synth → Real}) / median(d_{Real → Real})
% NNDR > 1 indicates synthetic records are farther from real than real
% records are from each other — the desired privacy regime.

\paragraph{Nearest-Neighbour Adversarial Accuracy (NNAA).}
\label{sec:eval:privacy:nnaa}
% Yale et al. \cite{healthgan} adversarial-accuracy test:
%   AA = 0.5 × (mean(d_{TS} > d_{TT}) + mean(d_{ST} > d_{SS}))

\paragraph{Membership Inference Attack.}
\label{sec:eval:privacy:mia}
% Threat model: adversary has access to synthetic data and a candidate
% record; must decide whether the record was in the training set.

\paragraph{Attribute Inference Attack.}
\label{sec:eval:privacy:aia}
% Threat model: adversary knows all but one sensitive attribute of a target
% record and tries to infer the missing attribute by querying the synthetic
% data.

\paragraph{Clustering Proximity / Outlier Linkability.}
\label{sec:eval:privacy:linkability}
% Identify outliers in real data: top 5 % of records ranked by distance from
% the centroid in StandardScaler feature space.

\subsubsection{Comparative Aggregation}
\label{sec:eval:aggregate}

% Maps to data_analysis.ipynb cells 49–52.

\paragraph{Privacy--Utility Tradeoff Scatter.}
\label{sec:eval:aggregate:tradeoff}
% Single 2D plot:
%   X = utility (Gradient-Boosting AUC-ROC under TSTR)
%   Y = privacy (NNDR median; higher = better)

\paragraph{Six-Dimension Radar Chart.}
\label{sec:eval:aggregate:radar}
% Six axes (all scaled to [0, 1] with higher = better).

\paragraph{Comprehensive Summary Table.}
\label{sec:eval:aggregate:table}
% 18 metrics × 5 sources (Real (TRTR) + 4 generators).

\cleardoublepage
